\documentclass[10pt,a4paper]{amsart}
\usepackage[margin=19mm]{geometry}
\usepackage{amsmath,amssymb,mathtools}
\usepackage[scr=rsfs,cal=euler]{mathalfa}
\usepackage{xfrac}
\usepackage[unicode,hidelinks,bookmarks=false]{hyperref}
\newcommand{\ie}{i.e.\ }
\newcommand{\cf}{cf.\ }
\newcommand{\resp}{respectively\ }
\newcommand{\Subsection}{\subsection}
\providecommand{\nobreakdash}{\nobreak\hbox{-}}
\setlength{\emergencystretch}{2em}
\multlinegap0pt
\usepackage{algorithm}\usepackage{algpseudocode}
\usepackage{amssymb}
\usepackage[shortlabels]{enumitem}
\usepackage{tikz}
\usepackage{xcolor}
\newtheorem{assumption}{Assumption}
\newtheorem{remark}{Remark}
\newtheorem{proposition}{Proposition}
\newtheorem{lemma}{Lemma}
\newtheorem{theorem}{Theorem}
\usepackage{cleveref}
\crefname{assumption}{Assumption}{Assumptions}
\crefname{remark}{Remark}{Remarks}
\crefname{proposition}{Proposition}{Propositions}
\crefname{lemma}{Lemma}{Lemmas}
\crefname{theorem}{Theorem}{Theorems}
\DeclareMathOperator{\Exp}{Exp}
\newcommand{\Hess}{\operatorname{Hess}}
\DeclareMathOperator{\grad}{grad}
\title{Nystr\"om Approximation on Manifolds}
\author{Hantao Nie, Bin Gao, Andi Han, Pratik Jawanpuria, Bamdev Mishra, Zaiwen Wen}
\date{Source: arXiv:2605.14933v1 (CC BY 4.0)}
\begin{document}
\maketitle

\noindent\emph{Source and licence.} Nystr\"om Approximation on Manifolds. Version arXiv:2605.14933v1 (CC BY 4.0), distributed by arXiv under the Creative Commons Attribution 4.0 International licence, http://creativecommons.org/licenses/by/4.0/, as recorded in the arXiv licence metadata for this exact version and shown on the abstract page https://arxiv.org/abs/2605.14933v1. Full licence notice: \texttt{licenses/arxiv-2605.14933-CC-BY-4.0.txt}.\par\smallskip

\noindent\textit{Editorial scope.} Counted unit: the article's theorem thm:riem-local-exp (source0.tex lines 2498--2520). The article's complete original proof of it is delivered with it (source0.tex lines 2522--2735, one \texttt{proof} environment of 214 lines). No mathematics is written, repaired or re-derived: the statement and the proof are the article's own lines, in the article's own notation, and no retained line is substituted or re-flowed.  Retained uncounted as the source-local context the proof needs: lines 1763--1766; lines 1776--1794; lines 1806--1818; lines 1844--1848; lines 1862--1869; lines 2069--2075; lines 2112--2121; lines 2299--2352; lines 2355--2367.  Nothing was omitted from any retained span, so the package carries no omission record.  No retained line contains a citation, so the wrapper carries no bibliography and no bibliography entry.  No retained line contains a figure environment, an image-inclusion command or drawing code, so no image is delivered and the package declares no asset dependency.  Editorial formatting only, in addition: an AMS \texttt{amsart} preamble that restates the article's own declarations and macros used by the retained lines, namely the environments \texttt{assumption}, \texttt{remark}, \texttt{proposition}, \texttt{lemma}, \texttt{theorem} on separate counters, together with the standard packages those lines need.  The retained environments print in this document as Assumption 1, Remark 1, Proposition 1, Lemma 1, Theorem 1 in order of appearance, whereas the article prints the same items with the numbering of its own full document.  The complete native source of the article as distributed in the arXiv e-print for this exact version is bundled unchanged as \texttt{sources/200624/source0.tex}.
\begin{equation}
\label{eq:cubic-model}
\phi(v)\;=\langle h_{x, B, \Xi}, v\rangle_x + \tfrac{1}{2}\langle \mathcal{J}_{x, B, \Xi} [v], v\rangle_x + \tfrac{\sigma}{6}\,\|v\|_x^{3},\qquad v\in B,
\end{equation}

\begin{algorithm}[H]
  \caption{Randomized Riemannian Nystr\"om cubic Newton method on $(\mathcal{M},g)$}
  \label{alg:sscn-riem}
  \begin{algorithmic}[1]
  \Require manifold $(\mathcal{M},g)$, objective function $f$, second-order retraction operator $R$, initial point $x_0\in\mathcal{M}$, sketch size $\ell < d$, cubic regularization parameter $\sigma_k>0$.
  \State Choose $x_0\in\mathcal{M}$.
  \For{$k=0,1,2,\dots, K-1$}
  \State Compute $\grad f(x_k)\in \mathrm{T}_{x_k}\mathcal{M}$.
  \State Choose subspace $B_k, \Xi_k$, map $\mathcal{F}_k$ satisfying the Haar--Grassmann sketching condition.
  \State Build $h_k:=\mathcal{P}_{x_k, B_k, \Xi_k}^* [\grad f(x_k)]$.
  \State Solve $v_k \in B_k$ from the subproblem:\[
  v_k \in \arg\min_{v\in B_k} \phi_k(v) := \langle h_k, v\rangle_{x_k} + \tfrac{1}{2}\langle \mathcal{J}_{x_k, B_k, \Xi_k} [v], v\rangle_{x_k} + \tfrac{\sigma_k}{6}\|v\|_{x_k}^{3}.
  \]
  \State Set $\eta_k:=\mathcal{P}_{x_k, B_k, \Xi_k}(v_k) \in \mathrm{T}_{x_k}\mathcal M$.
  \State Update $x_{k+1}=R_{x_k}(\eta_k)$.
  \State (Optional) Reject the update if $f(R_{x_k}(\eta_k)) > f(x_k)$.
  \EndFor
  \end{algorithmic}
  \end{algorithm}

\begin{assumption}
\label{ass:smooth}
The following conditions hold:
\begin{enumerate}
\item[1.] The function $f$ is twice continuously differentiable and geodesically convex on $\mathcal{M}$.

\item[2.] For every $x\in\mathcal{M}$, the pullback $\hat f:=f\circ R_x$ has an $L_3$-Lipschitz Hessian on $\mathrm{T}_x\mathcal{M}$, i.e., $\|\nabla^{2}\hat f(u)-\nabla^{2}\hat f(0)\|_{\mathrm{op}}\le L_3\|u\|_x$ for all $u\in\mathrm{T}_x\mathcal{M}$, and $\Hess f(x) \preceq L_2 \mathrm{Id}_x$.

\item[3.] The retraction operator $R_x$ is second-order on $\mathrm{T}_x\mathcal{M}$ for all $x\in\mathcal{M}$. 

\item[4.] Assume that the sequence $\{\mathcal{P}_k\}$ is independent and satisfies the Haar--Grassmann sketching condition.
\end{enumerate}
\end{assumption} 

\begin{assumption}
  (\textbf{strong geodesic convexity})
  \label{ass:strong-geo-convexity}
There exists a geodesically convex neighborhood $U$ of $x^\star$ such that $\Hess f(x)\succeq \mu \,\mathrm{Id}_x$ for all $x\in U$.
\end{assumption}

\begin{assumption}
  \label{ass:local-transported-resampling}
  Assume that $\mathcal P_{x^\star,B^\star,\Xi^\star}$ satisfies the Haar--Grassmann sketching condition at $x^\star$.
  Conditional on $x_k$, the operator $\widetilde{\mathcal P}_k$ is independent of the past and has the same distribution as 
  $
  \mathcal P_{x^\star,B^\star,\Xi^\star}
  $ at $x^\star$.
  \end{assumption}

  \begin{equation}
  \label{eq:retraction-cubic-upper-bound}
  \begin{aligned}
  f\big(R_x(\mathcal{P}_{x, B, \Xi} [v])\big) &\le f(x)+\langle h_{x, B, \Xi},v\rangle_x+\tfrac{1}{2}\langle \mathcal{J}_{x, B, \Xi} [v],v\rangle_x+\tfrac{M_\rho}{6}\|v\|_x^{3} ,\\
  f\big(R_x(\mathcal{P}_{x, B, \Xi} [v])\big) &\ge f(x)+\langle h_{x, B, \Xi},v\rangle_x+\tfrac{1}{2}\langle \mathcal{J}_{x, B, \Xi} [v],v\rangle_x - \tfrac{M_\rho}{6}\|v\|_x^{3}.
  \end{aligned}
  \end{equation}

  \begin{remark}
    \label{rem:cubic-upper-bound}
  The lemma shows that for $\rho$ sufficiently large, the solution to the subproblem~\eqref{eq:cubic-model} is in the interior of the ball of radius $\rho$. Without loss of generality, we can let $M = \max_{\rho \in [0, \bar \rho]} M_\rho$ be a uniform constant independent of $x, \rho$. Then~\eqref{eq:retraction-cubic-upper-bound} becomes
  \[
  \begin{aligned}
  f\big(R_x(\mathcal{P}_{x, B, \Xi} [v])\big) &\le f(x)+\langle h_{x, B, \Xi},v\rangle_x+\tfrac{1}{2}\langle \mathcal{J}_{x, B, \Xi} [v],v\rangle_x+\tfrac{M}{6}\|v\|_x^{3} ,\\
  f\big(R_x(\mathcal{P}_{x, B, \Xi} [v])\big) &\ge f(x)+\langle h_{x, B, \Xi},v\rangle_x+\tfrac{1}{2}\langle \mathcal{J}_{x, B, \Xi} [v],v\rangle_x - \tfrac{M}{6}\|v\|_x^{3}.
  \end{aligned}
  \]
  \end{remark}

\begin{proposition}
\label{prop:self-concordant}
Let $f$ be geodesically convex and $M_f$-self-concordant on an open geodesically convex neighborhood $U$.
\begin{enumerate}
\item[1.] For any $x,y\in U$, let
\[
v:=\Exp_x^{-1}(y)\in \mathrm{T}_x\mathcal{M},
\qquad
\delta_x(y):=\big\langle \Hess f(x)[v],v\big\rangle_x^{\frac{1}{2}}.
\]
If
$
M_f\,\delta_x(y)<1,
$
then
\[
\bigl(1-M_f\delta_x(y)\bigr)^2 \Hess f(x)
\preceq
\mathcal T_{y\to x}\,\Hess f(y)\,\mathcal T_{x\to y}
\preceq
\bigl(1-M_f\delta_x(y)\bigr)^{-2}\Hess f(x).
\]

\item[2.] For any $y\in U$, let
\[
\Lambda(y):=
\big\langle
\grad f(y),\,
[\Hess f(y)]^{-1}[\grad f(y)]
\big\rangle_y^{\frac{1}{2}}.
\]
If
$
M_f\,\Lambda(y)<1,
$
then
\[
f(y)-f(x^\star)\le \frac{1}{M_f^2}\,\omega_*\bigl(M_f\Lambda(y)\bigr),
\]
where
\[
\omega_*(t):=-t-\ln(1-t),
\quad t\in[0,1).
\]
In addition, for every $\gamma>0$ and every
$
t\in\Bigl[0,\frac{\gamma}{1+\gamma}\Bigr],
$
we have
$
\omega_*(t)\le \frac{1+\gamma}{2}\,t^2.
$
\end{enumerate}
\end{proposition}

\begin{lemma}
\label{lem:manifold-descent-lemma}
Let $\sigma_k\ge M$ in Algorithm~\ref{alg:sscn-riem}, where $M$ is defined in \cref{rem:cubic-upper-bound}. Then the exact cubic step satisfies
\[
f(x_k)-f(x_{k+1})
\ge
\frac{1}{2}\,
\Big\langle
h_k,\
\bigl(\mathcal J_k+\sqrt{\sigma_k/2}\,\|h_k\|_{x_k}^{\frac{1}{2}}\,\mathrm{Id}_{B_k}\bigr)^{-1}[h_k]
\Big\rangle_{x_k}.
\]
\end{lemma}

\begin{theorem}
\label{thm:riem-local-exp}
Suppose \cref{ass:smooth}, \cref{ass:strong-geo-convexity}, and \cref{ass:local-transported-resampling} hold. Let $\mathcal{H}^\star:=\Hess f(x^\star)$ and define
\[
\zeta
:=
\lambda_{\min}\!\Big(
(\mathcal{H}^\star)^{\frac{1}{2}}\,
\mathbb E\big[
\mathcal P_{x^\star,B^\star,\Xi^\star}
\bigl(\mathcal P_{x^\star,B^\star,\Xi^\star}^{*}\mathcal{H}^\star\mathcal P_{x^\star,B^\star,\Xi^\star}\bigr)^{-1}
\mathcal P_{x^\star,B^\star,\Xi^\star}^{*}
\big]\,
(\mathcal{H}^\star)^{\frac{1}{2}}
\Big).
\]
Then for any $\delta\in(0,1)$ there exists a neighborhood $V\subset U$ such that if $x_0\in V$, then
\[
\mathbb E\bigl[f(x_k)-f(x^\star)\bigr]
\le
\bigl(1-(1-\delta)\zeta\bigr)^k\,\bigl(f(x_0)-f(x^\star)\bigr).
\]
\end{theorem}

\begin{proof}
By isometric pullback, \cref{lem:manifold-descent-lemma} can be rewritten on the reference tangent space $\mathrm{T}_{x^\star}\mathcal{M}$ as
\begin{equation}
f(x_k)-f(x_{k+1})
\ge
\frac{1}{2}
\Big\langle
\widetilde h_k,\
\bigl(\widetilde{\mathcal J}_k+\tilde{\mu}_k\mathrm{Id}_{\widetilde B_k}\bigr)^{-1}[\widetilde h_k]
\Big\rangle_{x^\star},
\label{eq:descent-pulledback}
\end{equation}
where
$
\tilde{\mu}_k
:=
\sqrt{\sigma_k/2}\,\|\widetilde h_k\|_{x^\star}^{\frac{1}{2}}.
$
For any $\nu>0$, define
\[
S_{\nu}
:=
(\mathcal{H}^\star)^{\frac{1}{2}}\,
\mathbb E\Big[
\mathcal P_{x^\star,B^\star,\Xi^\star}
\bigl(
\mathcal P_{x^\star,B^\star,\Xi^\star}^{*}\mathcal{H}^\star\mathcal P_{x^\star,B^\star,\Xi^\star}
+\nu\,\mathrm{Id}_{B^\star}
\bigr)^{-1}
\mathcal P_{x^\star,B^\star,\Xi^\star}^{*}
\Big]\,
(\mathcal{H}^\star)^{\frac{1}{2}}.
\]
As $\nu\to 0$, the operators $S_{\nu}$ increase monotonically in the Loewner order to
\[
S_0
=
(\mathcal{H}^\star)^{\frac{1}{2}}\,
\mathbb E\Big[
\mathcal P_{x^\star,B^\star,\Xi^\star}
\bigl(
\mathcal P_{x^\star,B^\star,\Xi^\star}^{*}\mathcal{H}^\star\mathcal P_{x^\star,B^\star,\Xi^\star}
\bigr)^{-1}
\mathcal P_{x^\star,B^\star,\Xi^\star}^{*}
\Big]\,
(\mathcal{H}^\star)^{\frac{1}{2}},
\]
and therefore $\lambda_{\min}(S_{\nu})$ converges to $\lambda_{\min}(S_0)=\zeta$ as $\nu \to 0$.
Fix $\delta\in(0,1)$. Choose $\tau,\phi,\gamma>0$ such that
\begin{equation}
\frac{1-\tau}{(1+\phi)^2(1+\gamma)}\ge 1-\delta.
\label{eq:local-parameter-choice}
\end{equation}
Next, choose $\nu>0$ small enough such that
\begin{equation}
\lambda_{\min}(S_{\nu})\ge (1-\tau)\zeta.
\label{eq:Smu-close-to-zeta}
\end{equation}
By \cref{prop:self-concordant}.1, after shrinking $V\subset U$ if necessary, for all $x\in V$,
\begin{equation}
(1+\phi)^{-1}\mathcal{H}^\star
\preceq
\widetilde{\mathcal H}(x)
\preceq
(1+\phi)\mathcal{H}^\star.
\label{eq:Hessian-comparison-pulledback}
\end{equation}
In particular, for $x=x_k$, we have
\begin{equation}
\widetilde{\mathcal P}_k^*\,\widetilde{\mathcal H}_k\,\widetilde{\mathcal P}_k
\preceq
(1+\phi)\,\widetilde{\mathcal P}_k^*\mathcal{H}^\star\,\widetilde{\mathcal P}_k.
\label{eq:reduced-Hessian-comparison}
\end{equation}
Moreover, since $\widetilde g(x^\star)=0$ and $x\mapsto \widetilde g(x)$ is continuous, we have $\widetilde h_k\to 0$ as $x_k\to x^\star$. Hence, after shrinking $V$ once more if needed, we may assume that whenever $x_k\in V$,
\begin{equation}
\tilde{\mu}_k
=
\sqrt{\sigma_k/2}\,\|\widetilde h_k\|_{x^\star}^{\frac{1}{2}}
\le
(1+\phi)\nu.
\label{eq:augmentation-small}
\end{equation}
Combining~\eqref{eq:reduced-Hessian-comparison} and~\eqref{eq:augmentation-small}, we obtain
\[
\widetilde{\mathcal J}_k+\tilde{\mu}_k\mathrm{Id}_{\widetilde B_k}
\preceq
(1+\phi)\bigl(\widetilde{\mathcal P}_k^*\mathcal{H}^\star\,\widetilde{\mathcal P}_k+\nu\,\mathrm{Id}_{\widetilde B_k}\bigr),
\]
and therefore
\begin{equation}
\bigl(\widetilde{\mathcal J}_k+\tilde{\mu}_k\mathrm{Id}_{\widetilde B_k}\bigr)^{-1}
\succeq
\frac{1}{1+\phi}
\bigl(\widetilde{\mathcal P}_k^*\mathcal{H}^\star\,\widetilde{\mathcal P}_k+\nu\,\mathrm{Id}_{\widetilde B_k}\bigr)^{-1}.
\label{eq:inverse-comparison}
\end{equation}
Substituting~\eqref{eq:inverse-comparison} into~\eqref{eq:descent-pulledback} and using $\widetilde h_k=\widetilde{\mathcal P}_k^*[\widetilde g_k]$ yields
\[
f(x_k)-f(x_{k+1})
\ge
\frac{1}{2(1+\phi)}
\Big\langle
\widetilde g_k,\
\widetilde{\mathcal P}_k
\bigl(\widetilde{\mathcal P}_k^*\mathcal{H}^\star\,\widetilde{\mathcal P}_k+\nu\,\mathrm{Id}_{\widetilde B_k}\bigr)^{-1}
\widetilde{\mathcal P}_k^*[\widetilde g_k]
\Big\rangle_{x^\star}.
\]
Taking conditional expectation and using \cref{ass:local-transported-resampling}, we obtain
\begin{align}
&\quad \mathbb E\!\left[f(x_k)-f(x_{k+1})\,\middle|\,x_k\right] \\
&\ge
\frac{1}{2(1+\phi)}
\Big\langle
\widetilde g_k,\
\mathbb E\Big[
\mathcal P_{x^\star,B^\star,\Xi^\star}
\bigl(
\mathcal P_{x^\star,B^\star,\Xi^\star}^{*}\mathcal{H}^\star\mathcal P_{x^\star,B^\star,\Xi^\star}
+\nu\,\mathrm{Id}_{B^\star}
\bigr)^{-1}
\mathcal P_{x^\star,B^\star,\Xi^\star}^{*}
\Big][\widetilde g_k]
\Big\rangle_{x^\star} \nonumber\\
&=
\frac{1}{2(1+\phi)}
\Big\langle
(\mathcal{H}^\star)^{-1/2}\widetilde g_k,\
S_{\nu}\,(\mathcal{H}^\star)^{-1/2}\widetilde g_k
\Big\rangle_{x^\star} \nonumber\\
&\ge
\frac{\lambda_{\min}(S_{\nu})}{2(1+\phi)}
\big\langle
\widetilde g_k,\,
(\mathcal{H}^\star)^{-1}[\widetilde g_k]
\big\rangle_{x^\star}.
\label{eq:expected-descent-pulledback}
\end{align}
Define
$
\Lambda_k
:=
\big\langle
\widetilde g_k,\,
\widetilde{\mathcal H}_k^{-1}[\widetilde g_k]
\big\rangle_{x^\star}^{\frac{1}{2}}.
$
Since $\mathcal T_{x_k\to x^\star}$ is an isometry, it holds that
\[
\Lambda_k^2
=
\big\langle
\grad f(x_k),\,
[\Hess f(x_k)]^{-1}[\grad f(x_k)]
\big\rangle_{x_k}.
\]
By \cref{prop:self-concordant}.2, after shrinking $V$ if necessary, we may assume
$
M_f\Lambda_k\le \frac{\gamma}{1+\gamma}\quad \text{for all } x_k\in V.
$
Hence it follows that
\[
f(x_k)-f(x^\star)
\le
\frac{1}{M_f^2}\omega_*(M_f\Lambda_k)
\le
\frac{1+\gamma}{2}\Lambda_k^2
=
\frac{1+\gamma}{2}
\big\langle
\widetilde g_k,\,
\widetilde{\mathcal H}_k^{-1}[\widetilde g_k]
\big\rangle_{x^\star}.
\]
Using~\eqref{eq:Hessian-comparison-pulledback}, we have
\[
\widetilde{\mathcal H}_k^{-1}\preceq (1+\phi)(\mathcal{H}^\star)^{-1},
\]
and therefore
\begin{equation}
f(x_k)-f(x^\star)
\le
\frac{(1+\gamma)(1+\phi)}{2}
\big\langle
\widetilde g_k,\,
(\mathcal{H}^\star)^{-1}[\widetilde g_k]
\big\rangle_{x^\star}.
\label{eq:gap-upper-bound-pulledback}
\end{equation}
Combining~\eqref{eq:expected-descent-pulledback},~\eqref{eq:gap-upper-bound-pulledback}, and~\eqref{eq:Smu-close-to-zeta}, we get
\[
\mathbb E\!\left[f(x_k)-f(x_{k+1})\,\middle|\,x_k\right]
\ge
\frac{(1-\tau)\zeta}{(1+\phi)^2(1+\gamma)}\,\bigl(f(x_k)-f(x^\star)\bigr).
\]
Equivalently,
\[
\mathbb E\!\left[f(x_{k+1})-f(x^\star)\,\middle|\,x_k\right]
\le
\Bigl(1-\frac{(1-\tau)\zeta}{(1+\phi)^2(1+\gamma)}\Bigr)\,\bigl(f(x_k)-f(x^\star)\bigr).
\]
By~\eqref{eq:local-parameter-choice},
\[
\frac{(1-\tau)\zeta}{(1+\phi)^2(1+\gamma)}\ge (1-\delta)\zeta,
\]
and thus
\[
\mathbb E\!\left[f(x_{k+1})-f(x^\star)\,\middle|\,x_k\right]
\le
\bigl(1-(1-\delta)\zeta\bigr)\,\bigl(f(x_k)-f(x^\star)\bigr).
\]
Finally, by shrinking $V$ if necessary, the one-step decrease implies that the iterates remain in $V$, and the stated linear rate follows by iterating the inequality.
\end{proof}

\end{document}
